% ============================================================================
% Fichier  : partie2/P02_C05_risque_personnes.tex
% Titre    : Évaluer le Risque pour les Personnes
% Auteur   : MINKA MI NGUIDJOÏ Thierry Emmanuel
% Version  : 1.0 -- novembre 2026
% Corrigés : corriges/Corr_C05.tex
% ============================================================================
% Sources : loi n° 2024/017, art. 5 (évaluation d'impact), 25, 33 (texte
% officiel) ; ISO/IEC 29134 (étapes, paraphrasées) ; méthode PIA de la CNIL
% (échelles à quatre niveaux, paraphrasées) ; ISO 31000:2018.
% Faits nouveaux (bible) : évaluation d'impact du projet de téléconsultation
% du CENTRE MÉDICAL MBOLO, avril 2026.
% ============================================================================

\chapter{Évaluer le Risque pour les Personnes}
\label{chap:risque}
\chaptermark{Le risque pour les personnes}

\epigraph{%
    « Le risque que mesure la protection des données\\
    est celui que court la personne, pas l'entreprise. »%
}{Principe du cours}

\niveauChapitre{Intermédiaire}{%
    Chapitres~\ref{chap:principes} et~\ref{chap:cartographie}~; notions de
    gestion des risques (UE MSN-2S3-SEC08)
}

\begin{boxobjectifs}
\begin{itemize}
    \item Distinguer le risque pour les personnes du risque pour
          l'organisation.
    \item Appliquer les articles~25 et~33~: évaluation du risque pour tout
          traitement, évaluation d'impact pour les traitements à risque
          élevé.
    \item Repérer les traitements à risque élevé à l'aide de critères
          déclarés.
    \item Conduire une évaluation d'impact selon les étapes de
          l'ISO/IEC~29134~: description, nécessité, risques, mesures,
          décision.
    \item Auditer une évaluation d'impact existante.
\end{itemize}
\end{boxobjectifs}

\bigskip

% ============================================================================
\section*{Avant de commencer : « le risque, c'est la panne »}
% ============================================================================

Bertoua, lundi 6~avril~2026. Le CENTRE MÉDICAL MBOLO\footnote{Le CENTRE
MÉDICAL MBOLO, sa directrice, son personnel, ses patients et son prestataire
MEDISOFT SARL sont fictifs.} prépare un service de téléconsultation pour
les villages de la Kadey~: un infirmier équipé d'une tablette examine le
patient au centre de santé du village, et un médecin de Bertoua pose le
diagnostic par visioconférence. La directrice, le Dr~ABENA, a demandé au
responsable informatique, M.~YOUSSOUFA, une analyse des risques. Il revient
avec un tableau de quinze lignes~: coupure du réseau, panne de la tablette,
vol du matériel, perte de recettes, retard du prestataire.

Aucune ligne ne parle des patients~: de la consultation filmée qu'un tiers
pourrait voir, du diagnostic d'une maladie stigmatisante transmis par un
réseau partagé, de l'enfant examiné hors de la présence de ses parents. Le
tableau mesure ce que l'établissement peut perdre. La loi demande ce que
les personnes peuvent subir.

\begin{boxalert}
\textbf{Violation typique.} Présenter une analyse des risques pour
l'organisation (pannes, pertes financières) comme l'évaluation exigée par la
loi.

\smallskip
\textbf{Règle~:} l'évaluation d'impact apprécie « la probabilité et la
gravité des risques encourus par les droits et libertés des individus »
(article~5). Ses événements redoutés concernent les personnes~: la
divulgation, la modification, la disparition de leurs données, et leurs
conséquences sur leur vie.
\end{boxalert}

% ============================================================================
\section{Deux Obligations}
\label{sec:c05-obligations}
% ============================================================================

La loi impose deux évaluations distinctes.

\begin{boxjuris}
\textbf{\loicam{n°~2024/017}, article~25, alinéa~1}~\cite{loi2024017}~:
\textit{« Le responsable du traitement et le sous-traitant effectuent
l'évaluation du risque présenté par le traitement des données à caractère
personnel, suivant les critères d'évaluation, ainsi que les modalités
d'examen et de validation fixées par voie règlementaire. »}

\medskip
\textbf{\loicam{n°~2024/017}, article~33}~\cite{loi2024017}~:
\textit{« (1) Tout type de traitement susceptible d'engendrer un risque
élevé pour les droits et libertés des personnes physiques est subordonné à
la réalisation préalable d'une évaluation d'impact des opérations de
traitement envisagées sur la protection des données à caractère
personnel. (2) Les conditions et les modalités de la réalisation de
l'évaluation d'impact visées à l'alinéa 1 ci-dessus sont fixées par voie
règlementaire. »}

\smallskip
\textbf{Portée.} Tout traitement doit être évalué (article~25)~; seuls ceux
qui présentent un risque élevé exigent une évaluation d'impact préalable
(article~33). Dans les deux cas, critères et modalités relèvent de textes
d'application. En leur absence, l'auditeur vérifie que l'organisation a
déclaré et appliqué une méthode reconnue~; il cote \NEK{} de critère la
conformité aux modalités réglementaires elles-mêmes.
\end{boxjuris}

\begin{boxdef}
\termecle{Évaluation d'impact sur la protection des données}~: selon
l'article~5, « procédure visant à analyser la probabilité et la gravité des
risques encourus par les droits et libertés des individus résultant du
traitement de données à caractère personnel ».
\end{boxdef}

% ============================================================================
\section{Repérer le Risque Élevé}
\label{sec:c05-eleve}
% ============================================================================

La loi ne définit pas le « risque élevé ». En attendant le texte
d'application, le cours retient une liste de critères déclarés, inspirée des
pratiques internationales. Un traitement qui en réunit deux ou plus est
présumé à risque élevé et soumis à évaluation d'impact.

\begin{enumerate}
    \item Des données sensibles au sens de l'article~5 (santé, biométrie,
          transactions bancaires, origine régionale\ldots).
    \item Des personnes vulnérables~: mineurs, patients, salariés face à
          leur employeur, personnes en difficulté financière.
    \item Une décision automatisée ou un profilage produisant des effets sur
          la personne (Chapitre~\ref{chap:profilage}).
    \item Une surveillance systématique~: vidéosurveillance, géolocalisation,
          contrôle des communications.
    \item Un traitement à grande échelle, au regard du nombre de personnes
          ou de la durée.
    \item Le croisement ou l'interconnexion de fichiers (article~36).
    \item Une technologie nouvelle ou un usage nouveau d'une technologie.
    \item Un transfert hors du Cameroun (Chapitre~\ref{chap:transferts}).
\end{enumerate}

Ces critères sont une convention du cours, déclarée comme telle dans le
plan d'audit~; ils ne remplacent pas le texte d'application quand il sera
publié. La téléconsultation de MBOLO en réunit cinq~: données de santé,
patients vulnérables dont des mineurs, technologie nouvelle pour
l'établissement, transmission par un réseau d'opérateur, hébergement chez
un prestataire.

% ============================================================================
\section{Les Étapes d'une Évaluation d'Impact}
\label{sec:c05-etapes}
% ============================================================================

L'ISO/IEC~29134~\cite{iso29134} décrit le déroulement d'une évaluation
d'impact~; la méthode publiée par l'autorité française~\cite{cnil2018pia}
en donne une version outillée. On les résume ici en cinq étapes.

\begin{boxprocedure}
\textbf{Conduire une évaluation d'impact.}
\begin{enumerate}
    \item \textbf{Décrire} le traitement~: finalités, données, personnes,
          flux, supports, durées, acteurs.
    \item \textbf{Apprécier la nécessité et la proportionnalité}~: fondement,
          finalités, données strictement utiles, information, droits.
    \item \textbf{Apprécier les risques} pour les personnes~: sources de
          risque, événements redoutés, gravité, vraisemblance.
    \item \textbf{Choisir les mesures} qui ramènent chaque risque à un niveau
          acceptable, et évaluer le risque résiduel.
    \item \textbf{Décider et documenter}~: validation par le responsable,
          avis du délégué s'il existe, consultation des personnes ou de
          leurs représentants si possible, date de revue.
\end{enumerate}
\end{boxprocedure}

\subsection{Trois événements redoutés}

Les événements redoutés se rangent en trois familles, qui correspondent aux
trois propriétés de la sécurité vues du côté de la personne~:
\begin{itemize}
    \item l'\textbf{accès illégitime} à ses données (confidentialité)~: un
          voisin apprend un diagnostic~;
    \item la \textbf{modification non désirée} (intégrité)~: un résultat
          d'analyse attribué au mauvais patient~;
    \item la \textbf{disparition} (disponibilité)~: le dossier introuvable le
          jour d'une urgence.
\end{itemize}

\subsection{Gravité et vraisemblance}

\begin{tableaucours}{cL{3.0cm}L{8.0cm}}{%
    \tableauHeader{Niveau} &
    \tableauHeader{Gravité} &
    \tableauHeader{Effet sur la personne}
}
1 & Négligeable & Désagrément surmonté sans difficulté \\
2 & Limitée & Désagréments significatifs, surmontables avec quelques
efforts \\
3 & Importante & Conséquences importantes, surmontables avec de sérieuses
difficultés (exclusion, perte d'emploi, préjudice financier) \\
4 & Maximale & Conséquences irrémédiables (atteinte à la santé, à la
sécurité physique, stigmatisation durable) \\
\end{tableaucours}
\vspace{-0.8em}
\begin{center}{\small\itshape Tableau~5.1. Échelle de gravité (d'après la méthode de la CNIL, paraphrasée)}\end{center}

La vraisemblance se cote de même sur quatre niveaux (négligeable, limitée,
importante, maximale), d'après les sources de risque et les vulnérabilités
des supports. On reporte chaque risque sur une grille gravité ×
vraisemblance~; un risque dont la gravité ou la vraisemblance reste à~3
ou~4 après mesures n'est pas acceptable sans décision motivée.

\begin{boxinfo}
\textbf{Coter l'inconnu dans une évaluation d'impact.} Une vraisemblance
qu'on ne sait pas apprécier (par exemple, la sécurité réelle d'un
prestataire jamais audité) ne se cote pas « limitée » par défaut. On la
note \NEK{}, et l'on traite le risque comme s'il était au niveau le plus
élevé plausible tant que la preuve manque.
\end{boxinfo}

% ============================================================================
\section{La Téléconsultation de MBOLO}
\label{sec:c05-mbolo}
% ============================================================================

\begin{boxscenario}{Mission MBOLO. L'évaluation d'impact, avril 2026}
\textbf{Description.} Patients des villages de la Kadey, dont des enfants~;
données d'identité, symptômes, images de l'examen, diagnostic,
prescription~; flux tablette, réseau mobile, serveur de MEDISOFT à Douala,
poste du médecin à Bertoua~; enregistrement des séances envisagé « pour la
formation ».

\textbf{Nécessité.} Le fondement est la préservation de la santé
(article~9(2)) pour la consultation elle-même. L'enregistrement des séances
pour la formation n'a pas ce fondement~: il exige un consentement
spécifique, et des parents pour les mineurs (article~9(3)).

\textbf{Risques.} Trois risques prioritaires~:
\begin{itemize}
    \item R1~: un tiers présent dans la salle du centre de santé voit ou
          entend la consultation (gravité~3, vraisemblance~3)~;
    \item R2~: interception ou accès indu aux images stockées chez le
          prestataire (gravité~4, vraisemblance \NEK{}, prestataire jamais
          audité)~;
    \item R3~: erreur d'attribution d'une image au mauvais patient
          (gravité~4, vraisemblance~2).
\end{itemize}

\textbf{Mesures.} Salle fermée et procédure d'accueil~; chiffrement de bout
en bout et contrat avec MEDISOFT révisé (Chapitre~\ref{chap:soustraitants})~;
double identification du patient~; abandon de l'enregistrement des
séances.

\textbf{Décision} (8~mai~2026). R1 et R3 ramenés à un niveau acceptable.
R2 reste \NEK{} tant que MEDISOFT n'a pas été audité~: le Dr~ABENA lance le
service sur deux villages seulement, en attendant l'audit du prestataire.
\end{boxscenario}

\begin{boxlecture}{L'évaluation d'impact de MBOLO, auditée en octobre 2026}
\begin{tableaucours}{L{7.8cm}cL{3.2cm}}{%
    \tableauHeader{Exigence} &
    \tableauHeader{Cote} &
    \tableauHeader{Preuve}
}
L'évaluation a précédé la mise en service & \OUI{} & Datée du 8~mai \\
Le traitement est décrit avec ses flux & \OUI{} & Schéma, liste des
données \\
Le fondement de chaque finalité est apprécié & \OUI{} & Section
nécessité \\
Les risques sont appréciés du point de vue des patients & \OUI{} & Trois
risques, échelles \\
Les mesures décidées sont en place & \NON{} & Chiffrement de bout en bout
non activé \\
Le risque lié au prestataire est apprécié & \NEK{} & Audit de MEDISOFT
non fait \\
La conformité aux modalités réglementaires est établie & \NEK{} & Texte
d'application non publié \\
\end{tableaucours}

\medskip
\textbf{Décompte.} Quatre \OUI{}, un \NON{}, deux \NEK{} sur sept.
\hfill\lecture{0{,}57}{0{,}29}{1{,}38}

\medskip
\textbf{Lecture.} L'évaluation est bien faite sur le papier, mais
$U = 0{,}29$~: le domaine reste opaque. Les deux \NEK{} sont de nature
différente~: l'un est de preuve (l'audit du prestataire le lèvera), l'autre
de critère (seul le texte d'application le lèvera). Le \NON{} est un écart
entre la décision et l'exécution, le plus fréquent dans les évaluations
d'impact.
\end{boxlecture}

% ============================================================================
\section{Auditer une Évaluation d'Impact}
\label{sec:c05-auditer}
% ============================================================================

L'auditeur ne refait pas l'évaluation~; il vérifie qu'elle existe quand elle
est exigée, qu'elle a été faite \textit{avant} le traitement, qu'elle
porte sur les personnes, que ses cotes sont justifiées, et surtout que les
mesures décidées sont en place. Ce dernier point est le plus souvent en
défaut~: une évaluation d'impact décrit un traitement idéal, et le
traitement réel s'en écarte.

Pour les traitements qui n'en ont pas, l'auditeur applique les critères de
la section~\ref{sec:c05-eleve} aux fiches du registre. À NYANGA CRÉDIT,
douze des quarante-sept traitements réunissent au moins deux critères~;
aucun n'a fait l'objet d'une évaluation d'impact. Le classeur des mauvais
payeurs en réunit quatre.

% ============================================================================
\section{Ce que le Risque demande à chaque Métier}
\label{sec:c05-metiers}
% ============================================================================

\begin{boxmetier}{Responsable de traitement}
Décider, au vu de l'évaluation, de lancer, d'adapter ou de renoncer, et
signer cette décision. Le Dr~ABENA a limité le service à deux villages~:
c'est une décision proportionnée à un risque encore inconnu.
\end{boxmetier}

\begin{boxmetier}{Délégué à la protection des données}
Conduire ou animer l'évaluation d'impact, rendre un avis écrit, et suivre la
mise en place des mesures jusqu'à la preuve.
\end{boxmetier}

\begin{boxmetier}{RSSI}
Apporter l'appréciation de la vraisemblance~: vulnérabilités des supports,
menaces réelles, état des mesures. Signaler quand une mesure décidée n'est
pas activée.
\end{boxmetier}

\begin{boxresume}
\begin{itemize}
    \item La loi exige une évaluation du risque pour tout traitement
          (article~25) et une évaluation d'impact préalable pour les
          traitements à risque élevé (article~33).
    \item Le risque évalué est celui des personnes~: accès illégitime,
          modification non désirée, disparition de leurs données.
    \item En l'absence de texte d'application, le risque élevé se repère à
          l'aide de critères déclarés, au moins deux réunis.
    \item Une évaluation d'impact se juge surtout à la mise en place
          effective des mesures qu'elle décide.
    \item À MBOLO, l'évaluation est sérieuse mais une mesure n'est pas
          activée et le risque lié au prestataire reste \NEK{}.
\end{itemize}
\end{boxresume}

\begin{boxinfo}
\textbf{Points de contrôle.} Vous avez compris ce chapitre si vous savez~:
\begin{itemize}
    \item distinguer les articles~25 et~33~;
    \item dire si un traitement exige une évaluation d'impact, critères à
          l'appui~;
    \item coter la gravité et la vraisemblance d'un risque pour une
          personne.
\end{itemize}
\end{boxinfo}

\begin{boxexo}[title={\ding{45}\quad Exercice~5.1. Évaluation d'impact ou non \niveaubase}]
Dites si ces traitements exigent une évaluation d'impact, critères à
l'appui~: (a) la paie de 210 salariés~; (b) le pointage biométrique du
personnel~; (c) la liste des fournisseurs~; (d) la notation automatique du
crédit chez NOUN PAY~; (e) la vidéosurveillance d'une salle d'attente de
pédiatrie.
\end{boxexo}

\begin{boxexo}[title={\ding{45}\quad Exercice~5.2. Réécrire le tableau \niveaumoyen}]
Reprenez le tableau de M.~YOUSSOUFA. Pour trois de ses lignes, écrivez
l'événement redouté du point de vue du patient, sa gravité et sa
vraisemblance.
\end{boxexo}

\begin{boxexo}[title={\ding{45}\quad Exercice~5.3. Le classeur des mauvais payeurs \niveaumoyen}]
Conduisez les trois premières étapes de l'évaluation d'impact du classeur
des mauvais payeurs de NYANGA. Quelle conclusion tirez-vous de l'étape de
nécessité~?
\end{boxexo}

\begin{boxexo}[title={\ding{45}\quad Exercice~5.4. Un risque inconnu \niveauexpert}]
Le Dr~ABENA vous demande s'il peut étendre la téléconsultation à vingt
villages avant l'audit de MEDISOFT. Rédigez votre avis, en séparant les
constatations de l'analyse, et en indiquant ce qui ferait passer R2 de
\NEK{} à une cote établie.
\end{boxexo}

\bigskip

\begin{boxinfo}
\textbf{Pour aller plus loin.}
\begin{itemize}
    \item L'ISO/IEC~29134~\cite{iso29134}, pour les étapes de l'évaluation
          d'impact.
    \item La méthode PIA de la CNIL~\cite{cnil2018pia}, librement
          accessible, pour ses échelles et ses modèles.
\end{itemize}
\end{boxinfo}

\begin{boxinfo}
\textbf{Transition.} Parmi les risques pour les personnes, l'un tient à leur
capacité d'agir sur leurs propres données. Peuvent-elles savoir ce que l'on
sait d'elles, le corriger, s'y opposer~? C'est l'objet du
Chapitre~\ref{chap:droits}.
\end{boxinfo}
